\documentclass[12pt]{article}
\usepackage[utf8]{inputenc}
\usepackage[T1]{fontenc}
\usepackage{hyperref}
\title{Manangkari Language Resources for K-12 Curriculum Development}
\author{Compiled by Claude Sonnet 5.5 (claude-sonnet-5-5) for Andr\'as Kornai}
\date{}
\begin{document}
\maketitle

\section{Introduction}
Manangkari (ISO 639-3 znk; also spelled Mananggari and Naragani) is an extinct Australian language. Search results place it on the Goulburn Islands, probably mainly North Goulburn Island, in the Northern Territory. They describe it as probably a dialect of, or closely related to, Maung (Mawng), an Iwaidjan language that is still spoken. No native speakers are known. Very little material appears to exist, and none of it is designed for schools.

\section{Primary Educational Resources}
No primers, textbooks or school programmes for Manangkari were found in search results. Any heritage curriculum would have to draw on the closely related Maung language, and that material should not be presented as Manangkari.

\section{Linguistic and Reference Materials}
\begin{itemize}
\item Wikipedia, Manangkari language (a short overview): \url{https://en.wikipedia.org/wiki/Manangkari_language}
\item AIATSIS Austlang record for Mananggari (N183), with the language's location and its relationship to Mawng: \url{https://collection.aiatsis.gov.au/austlang/language/N183}
\item Endangered Languages Project entry: \url{https://www.endangeredlanguages.com/lang/2546}
\item Wikipedia, Maung language (the related living language): \url{https://en.wikipedia.org/wiki/Maung_language}
\end{itemize}

\section{Religious and Community Resources}
None were found.

\section{Digital and Audio Resources}
No audio, video or corpus for Manangkari was found.

\section{Curriculum Development Context}
The language has no known speakers and almost no documentation. A curriculum could therefore only be a heritage awareness programme, covering the history of the Goulburn Islands and the place of Manangkari within the Iwaidjan languages. Any language content would depend on archival research, for example through AIATSIS, and on consultation with the Aboriginal communities concerned. Those communities, probably Maung-speaking ones, would need to decide what is appropriate.

\end{document}